\section{Tareas en GPU}

La implementación de los filtros Sobel y Laplace es directamente paralelizable procesando cada píxel por separado. Para Fourier y Canny utilizamos librerías de NVIDIA, que ya están optimizadas para GPU.

Para el renderizado, trabajamos con una técnica estándar en este tipo de algoritmo que es el basado en tiles (o mosaicos). En esta técnica, la imagen de salida se divide en regiones de tamaño fijo, en este caso $16 \times 16$ píxeles, y cada región es procesada por un bloque de hilos en la GPU. Cada hilo dentro del bloque se encarga de calcular la contribución de las gaussianas a un píxel específico dentro de esa región. El motivo de emplear esta división es que las gaussianas contribuyen de forma significativa a un área limitada de la imagen, por lo que los píxeles de una misma región comparten muchas de las gaussianas, mientras que para regiones separadas las gaussianas relevantes varían. Así, podemos descartar muchas gaussianas según su distancia a la región, reduciendo los cálculos necesarios por píxel y aprovechando la memoria compartida del bloque para las gaussianas involucradas.

Para las jacobianas, y el resto de operaciones relacionadas, trabajamos con un hilo por gaussiana, donde el resultado numérico se acumula en la memoria global antes de ser pasado a la CPU. De este modo, y al igual que en VIGS-Fusion, la CPU solo trabaja con matrices de tamaño reducido.



